% Formulario TEEA 3MAT — Argomento 3: resistività, temperatura, potenza ed energia
% Per rigenerare il PDF (dalla cartella strumenti, due volte per il numero di pagine):
%   pdflatex -output-directory=../pdf -jobname=03-formulario-resistivita-potenza-energia formulario.tex
% ---------------------------------------------------------------------------
%  Stile comune dei formulari TEEA (pdfLaTeX, solo pacchetti standard
%  presenti in TeX Live e MiKTeX). Testo in Computer Modern Sans.
% ---------------------------------------------------------------------------
\documentclass[11pt,a4paper]{article}
\usepackage[a4paper, top=14mm, bottom=16mm, left=14mm, right=14mm, headheight=0pt, footskip=8mm]{geometry}
\usepackage[utf8]{inputenc}
\usepackage{amsmath, amssymb}
\renewcommand{\familydefault}{\sfdefault}
\newcommand{\ohm}{\ensuremath{\Omega}}
\newcommand{\uA}{\ensuremath{\mu}}
\newcommand{\per}{\ensuremath{\cdot}}
\usepackage{xcolor, tikz, multicol, booktabs, array, fancyhdr, lastpage, tabularx}
\usepackage[most]{tcolorbox}
\usepackage[siunitx, american resistors, american currents]{circuitikz}
\ctikzset{bipoles/thickness=1, current/distance=0.35}
\usetikzlibrary{arrows.meta}

\definecolor{verde}{HTML}{159957}
\definecolor{blu}{HTML}{155799}
\definecolor{inchiostro}{HTML}{1B2936}
\definecolor{tenue}{HTML}{606C71}
\definecolor{fondo}{HTML}{EEF3F0}
\definecolor{tensione}{HTML}{2152C0}
\definecolor{corrente}{HTML}{B3205F}
\definecolor{resistenza}{HTML}{0B7A62}
\color{inchiostro}

\setlength{\parindent}{0pt}
\setlength{\columnsep}{7mm}
\renewcommand{\arraystretch}{1.25}

% unità di misura in tondo e virgola decimale
\newcommand{\U}[1]{\,\mathrm{#1}}
\newcommand{\virg}{{,}}

% riquadro di un argomento
\newtcolorbox{blocco}[1]{enhanced, breakable=false, colback=white, colframe=verde, boxrule=0.6pt,
  leftrule=3pt, arc=1.5pt, left=6pt, right=6pt, top=4pt, bottom=5pt,
  title={\bfseries #1}, coltitle=verde, colbacktitle=fondo, fonttitle=\normalsize,
  toptitle=1pt, bottomtitle=1pt, before skip=0pt, after skip=4.5mm}

% formula grande centrata
\newcommand{\F}[1]{\par\smallskip{\centering\Large$\displaystyle #1$\par}\smallskip}
% legenda dei simboli
\newcommand{\leg}[1]{{\small\color{tenue}#1}\par}

% intestazione del formulario
\newcommand{\intestazione}[3]{%
  {\color{verde}\bfseries\small #1}\par\vspace{1mm}
  {\fontsize{19}{23}\selectfont\bfseries #2}\par\vspace{1mm}
  {\small\color{tenue} #3}\par\vspace{2mm}
  {\color{verde}\rule{\linewidth}{1.6pt}}\par\vspace{4mm}}

\pagestyle{fancy}
\fancyhf{}
\renewcommand{\headrulewidth}{0pt}
\fancyfoot[L]{\scriptsize\color{tenue}Prof. Pasquale Naclerio \per{} TEEA 3MAT \per{} Formulario}
\fancyfoot[R]{\scriptsize\color{tenue}Pagina \thepage\ di \pageref{LastPage}}

\begin{document}

\intestazione{TEEA --- 3MAT \per{} Argomento 3 \per{} FORMULARIO}
             {Resistività, temperatura, potenza ed energia}
             {Prof. Pasquale Naclerio \per{} Manutenzione e Assistenza Tecnica}

\begin{multicols}{2}

\begin{blocco}{Resistenza di un conduttore}
\F{R = \rho \cdot \frac{L}{S}}
\leg{$R$ [\ohm] \quad $\rho$ [\ohm$\cdot$mm$^2$/m] \quad $L$ [m] \quad $S$ [mm$^2$]}
\textbf{filo tondo:}
\F{S = \frac{\pi \cdot d^2}{4} \qquad d = \sqrt{\frac{4 \cdot S}{\pi}}}
\textbf{pista (sezione rettangolare):}
\F{S = w \cdot t}
\leg{$w$ larghezza [mm] \quad $t$ spessore [mm]}
\leg{spessore tipico del rame: $t = 35\U{\mu m} = 0\virg035\U{mm}$}
\end{blocco}

\begin{blocco}{Stessa resistenza con due materiali}
\F{\frac{S_{Al}}{S_{Cu}} = \frac{\rho_{Al}}{\rho_{Cu}} = 1\virg6}
\leg{densità: rame $8\virg96\U{g/cm^3}$ \quad alluminio $2\virg70\U{g/cm^3}$}
\end{blocco}

\begin{blocco}{Resistenza e temperatura}
\F{R(T) = R_0 \cdot \left[1 + \alpha \cdot (T - T_0)\right]}
\F{\Delta T = T - T_0}
\leg{$R_0$ resistenza a $T_0$ (di solito $20\U{^\circ C}$) \quad $\alpha$ [$1/^\circ$C]}
\end{blocco}

\begin{blocco}{Materiali (a $20\U{^\circ C}$)}
\begin{tabularx}{\linewidth}{@{}X r r@{}}
\toprule
 & $\rho$ [\ohm$\cdot$mm$^2$/m] & $\alpha$ [$1/^\circ$C]\\
\midrule
rame (Cu) & $0\virg0175$ & $0\virg0039$\\
alluminio (Al) & $0\virg028$ & $0\virg0040$\\
tungsteno (W) & $0\virg055$ & $0\virg0045$\\
costantana & $0\virg49$ & $\approx 0$\\
nichel-cromo & $1\virg1$ & ---\\
\bottomrule
\end{tabularx}
\end{blocco}

\columnbreak

\begin{blocco}{Potenza ed effetto Joule}
\F{P = V \cdot I \qquad P = R \cdot I^2 \qquad P = \frac{V^2}{R}}
\F{Q = P \cdot \Delta t = R \cdot I^2 \cdot \Delta t}
\leg{$I$ imposta: $R\uparrow \;\Rightarrow\; P\uparrow$ \qquad $V$ imposta: $R\uparrow \;\Rightarrow\; P\downarrow$}
\end{blocco}

\begin{blocco}{Energia e costo}
\F{E = P \cdot t}
\leg{$E$ [Wh] \quad $P$ [W] \quad $t$ [h]}
\F{1\U{kWh} = 1000\U{Wh}}
\F{1\U{Wh} = 3600\U{J} \qquad 1\U{J} = 1\U{W}\cdot 1\U{s}}
\leg{costo [euro] $=$ $E$ [kWh] $\cdot$ prezzo [euro/kWh]}
\end{blocco}

\begin{blocco}{Linea e carico in serie (partitore di tensione)}
\begin{minipage}[c]{0.34\linewidth}\centering
\begin{circuitikz}[scale=0.8, transform shape, line width=0.7pt]
  \draw[line width=1.1pt] (-0.6,4.4) -- (0.6,4.4); \node[above] at (0,4.4) {$V$};
  \draw (0,4.4) to[R, l_=$R_L$, i>^=$I$] (0,2.2) to[R, l_=$R_C$, *-] (0,0) node[sground]{};
  \draw (0,2.2) to[short, -o] (1.1,2.2) node[right]{$X$};
\end{circuitikz}
\end{minipage}\hfill
\begin{minipage}[c]{0.64\linewidth}
\F{V_C = V \cdot \frac{R_C}{R_L + R_C}}
\F{V_L = V \cdot \frac{R_L}{R_L + R_C}}
\F{\eta = \frac{P_C}{P_{tot}} = \frac{V_C}{V}}
\end{minipage}
\end{blocco}

\begin{blocco}{Carichi in parallelo (partitore di corrente)}
\F{I_1 = I \cdot \frac{R_2}{R_1 + R_2} \qquad I_2 = I \cdot \frac{R_1}{R_1 + R_2}}
\F{P_{tot} = P_1 + P_2 \qquad E_{tot} = E_1 + E_2}
\end{blocco}

\end{multicols}

\end{document}
